% P10-4b, 14.09.2026: Begriffs-/Beleganschlüsse und abgegrenzte Prüfaussagen präzisiert.
% M08, 11.09.2026: Gestaltungsanlässe erhalten; Fallzahlen E48-01 im Anhang.
\section{Projektionen und kontrollierte Außenkopplung}
\label{sec:projektionen-aussenkopplung}

\emph{Zielbild und Designproblem.} Die Übergabe nach GitHub gehörte
zum ursprünglichen Zielkorridor. Mit dem persistenten Projektzustand
aus Abschnitt~\ref{sec:projektzustand-kontrollierte-mutation} entstand
eine neue Frage: Wie können Issues und Dokumente operativ genutzt
und verändert werden, ohne einen konkurrierenden fachlichen Bestand
zu führen? Die Entwurfsentscheidung unterscheidet deshalb die
maßgebliche Projektinformation von ihren abgeleiteten Außensichten.
Externe Änderungen sollen als neue Eingaben beurteilt werden und
den Bestand nicht still überschreiben.

\emph{Formative Befunde und Entwurfsfolge.} Ein dokumentierter
Durchlauf zeigte, dass Arbeitselemente als Issues angelegt und nach
fachlichen Änderungen gezielt aktualisiert werden konnten (E48-01).
Die Zuordnung wurde auf höchstens ein Issue je Arbeitselement und
höchstens ein Arbeitselement je Issue begrenzt. Ein weiterer Befund
zeigte jedoch, dass eine korrekte Zuordnungsregel einen falschen
Außenbestand nicht kompensiert: Ein Synchronisationslauf verwendete
einen veralteten Schnappschuss eines anderen Repositories und hätte
22 falsche Verknüpfungen erzeugt. Die vorgelagerte Probelauf-Stufe
machte dies vor dem Schreiben sichtbar (E48-02). Erhalten ist der
zeitnahe Protokolleintrag; der damalige Laufordner wurde entfernt.
Daraus folgten die Neuerhebung des Außenbestands als Ablaufprinzip
und eine technische Prüfung seiner Zugehörigkeit zum Ziel-Repository.

Auch die Inhaltsbildung wurde eingegrenzt: Ein modellformulierter
Aktualisierungsvorschlag hatte einen ausführlichen Issue-Text durch
eine dünne Statuszeile ersetzt. Daraufhin wurden die maßgeblichen
Textbausteine für Anlage und Aktualisierung in einem gemeinsamen
deterministischen Erzeugungsvertrag zusammengeführt. Er legt fest,
welche Core-Inhalte in welcher Form dargestellt werden.

Für die Gegenrichtung dokumentiert eine spätere Abnahmerunde einen
externen Issue-Kommentar, der mit Herkunftsangabe als Vorschlag
aufbereitet und dem Einordnungs- und Autorisierungsweg zugeführt
wurde. Die Entscheidung wurde anschließend als Vermerk am Issue
sichtbar. Im Folgeabruf wurde der bereits verarbeitete Kommentar
ausgefiltert; andere Funde blieben bestehen (E48-03). Hin- und Rückweg übernehmen damit unterschiedliche
Aufgaben: Der eine stellt bereits geltende Inhalte dar, der andere
bereitet externe Informationen für eine erneute fachliche Beurteilung
vor.

Das Prinzip gilt auch für maschinell gepflegte Dokument- und
Architektursichten; redaktionelle Bereiche bleiben davon getrennt.
Rückwege bestehen nur für die vorgesehenen operativen Kanäle.
Daraus ergibt sich die achte Designanforderung:

\textbf{A8:} Operative Darstellungen und externe Systeme dürfen keinen
konkurrierenden fachlichen Projektbestand bilden; ihre maßgeblichen
Inhalte müssen aus dem autorisierten Projektzustand abgeleitet und
eindeutig auf ihn bezogen werden. Externe Rückmeldungen sind als
kontrollierte Eingänge zu verarbeiten; daraus resultierende
Zustandsänderungen dürfen nur über die bestehenden Mutations- und
Autorisierungspfade wirksam werden.

\emph{Einordnung und Grenzen.} Die Fälle belegen ausgeführte
Projektions- und Rückwegschritte sowie die Wirkung des Probelaufs
im bezeichneten Fehlerfall. Sie belegen weder die vollständige
Erfassung externer Änderungen noch allgemeine Drift-Freiheit.
Abschnitt~\ref{sec:projektionen-kontrollierte-aussenkopplung} erklärt
die resultierende Architektur; Kapitel~\ref{chap:evaluation}
untersucht ausgewählte Übergaben und Kontrollfälle.
